% Options for packages loaded elsewhere
\PassOptionsToPackage{unicode}{hyperref}
\PassOptionsToPackage{hyphens}{url}
%
\documentclass[
]{article}
\usepackage{amsmath,amssymb}
\usepackage{iftex}
\ifPDFTeX
  \usepackage[T1]{fontenc}
  \usepackage[utf8]{inputenc}
  \usepackage{textcomp} % provide euro and other symbols
\else % if luatex or xetex
  \usepackage{unicode-math} % this also loads fontspec
  \defaultfontfeatures{Scale=MatchLowercase}
  \defaultfontfeatures[\rmfamily]{Ligatures=TeX,Scale=1}
\fi
\usepackage{lmodern}
\ifPDFTeX\else
  % xetex/luatex font selection
\fi
% Use upquote if available, for straight quotes in verbatim environments
\IfFileExists{upquote.sty}{\usepackage{upquote}}{}
\IfFileExists{microtype.sty}{% use microtype if available
  \usepackage[]{microtype}
  \UseMicrotypeSet[protrusion]{basicmath} % disable protrusion for tt fonts
}{}
\makeatletter
\@ifundefined{KOMAClassName}{% if non-KOMA class
  \IfFileExists{parskip.sty}{%
    \usepackage{parskip}
  }{% else
    \setlength{\parindent}{0pt}
    \setlength{\parskip}{6pt plus 2pt minus 1pt}}
}{% if KOMA class
  \KOMAoptions{parskip=half}}
\makeatother
\usepackage{xcolor}
\usepackage[margin=1in]{geometry}
\usepackage{color}
\usepackage{fancyvrb}
\newcommand{\VerbBar}{|}
\newcommand{\VERB}{\Verb[commandchars=\\\{\}]}
\DefineVerbatimEnvironment{Highlighting}{Verbatim}{commandchars=\\\{\}}
% Add ',fontsize=\small' for more characters per line
\usepackage{framed}
\definecolor{shadecolor}{RGB}{248,248,248}
\newenvironment{Shaded}{\begin{snugshade}}{\end{snugshade}}
\newcommand{\AlertTok}[1]{\textcolor[rgb]{0.94,0.16,0.16}{#1}}
\newcommand{\AnnotationTok}[1]{\textcolor[rgb]{0.56,0.35,0.01}{\textbf{\textit{#1}}}}
\newcommand{\AttributeTok}[1]{\textcolor[rgb]{0.13,0.29,0.53}{#1}}
\newcommand{\BaseNTok}[1]{\textcolor[rgb]{0.00,0.00,0.81}{#1}}
\newcommand{\BuiltInTok}[1]{#1}
\newcommand{\CharTok}[1]{\textcolor[rgb]{0.31,0.60,0.02}{#1}}
\newcommand{\CommentTok}[1]{\textcolor[rgb]{0.56,0.35,0.01}{\textit{#1}}}
\newcommand{\CommentVarTok}[1]{\textcolor[rgb]{0.56,0.35,0.01}{\textbf{\textit{#1}}}}
\newcommand{\ConstantTok}[1]{\textcolor[rgb]{0.56,0.35,0.01}{#1}}
\newcommand{\ControlFlowTok}[1]{\textcolor[rgb]{0.13,0.29,0.53}{\textbf{#1}}}
\newcommand{\DataTypeTok}[1]{\textcolor[rgb]{0.13,0.29,0.53}{#1}}
\newcommand{\DecValTok}[1]{\textcolor[rgb]{0.00,0.00,0.81}{#1}}
\newcommand{\DocumentationTok}[1]{\textcolor[rgb]{0.56,0.35,0.01}{\textbf{\textit{#1}}}}
\newcommand{\ErrorTok}[1]{\textcolor[rgb]{0.64,0.00,0.00}{\textbf{#1}}}
\newcommand{\ExtensionTok}[1]{#1}
\newcommand{\FloatTok}[1]{\textcolor[rgb]{0.00,0.00,0.81}{#1}}
\newcommand{\FunctionTok}[1]{\textcolor[rgb]{0.13,0.29,0.53}{\textbf{#1}}}
\newcommand{\ImportTok}[1]{#1}
\newcommand{\InformationTok}[1]{\textcolor[rgb]{0.56,0.35,0.01}{\textbf{\textit{#1}}}}
\newcommand{\KeywordTok}[1]{\textcolor[rgb]{0.13,0.29,0.53}{\textbf{#1}}}
\newcommand{\NormalTok}[1]{#1}
\newcommand{\OperatorTok}[1]{\textcolor[rgb]{0.81,0.36,0.00}{\textbf{#1}}}
\newcommand{\OtherTok}[1]{\textcolor[rgb]{0.56,0.35,0.01}{#1}}
\newcommand{\PreprocessorTok}[1]{\textcolor[rgb]{0.56,0.35,0.01}{\textit{#1}}}
\newcommand{\RegionMarkerTok}[1]{#1}
\newcommand{\SpecialCharTok}[1]{\textcolor[rgb]{0.81,0.36,0.00}{\textbf{#1}}}
\newcommand{\SpecialStringTok}[1]{\textcolor[rgb]{0.31,0.60,0.02}{#1}}
\newcommand{\StringTok}[1]{\textcolor[rgb]{0.31,0.60,0.02}{#1}}
\newcommand{\VariableTok}[1]{\textcolor[rgb]{0.00,0.00,0.00}{#1}}
\newcommand{\VerbatimStringTok}[1]{\textcolor[rgb]{0.31,0.60,0.02}{#1}}
\newcommand{\WarningTok}[1]{\textcolor[rgb]{0.56,0.35,0.01}{\textbf{\textit{#1}}}}
\usepackage{graphicx}
\makeatletter
\def\maxwidth{\ifdim\Gin@nat@width>\linewidth\linewidth\else\Gin@nat@width\fi}
\def\maxheight{\ifdim\Gin@nat@height>\textheight\textheight\else\Gin@nat@height\fi}
\makeatother
% Scale images if necessary, so that they will not overflow the page
% margins by default, and it is still possible to overwrite the defaults
% using explicit options in \includegraphics[width, height, ...]{}
\setkeys{Gin}{width=\maxwidth,height=\maxheight,keepaspectratio}
% Set default figure placement to htbp
\makeatletter
\def\fps@figure{htbp}
\makeatother
\setlength{\emergencystretch}{3em} % prevent overfull lines
\providecommand{\tightlist}{%
  \setlength{\itemsep}{0pt}\setlength{\parskip}{0pt}}
\setcounter{secnumdepth}{-\maxdimen} % remove section numbering
\ifLuaTeX
  \usepackage{selnolig}  % disable illegal ligatures
\fi
\IfFileExists{bookmark.sty}{\usepackage{bookmark}}{\usepackage{hyperref}}
\IfFileExists{xurl.sty}{\usepackage{xurl}}{} % add URL line breaks if available
\urlstyle{same}
\hypersetup{
  pdftitle={Bericht 1 Fallstudien},
  hidelinks,
  pdfcreator={LaTeX via pandoc}}

\title{Bericht 1 Fallstudien}
\author{}
\date{\vspace{-2.5em}2025-10-17}

\begin{document}
\maketitle

Die erhobenen Daten sind Erhebungen zu dem Thema Glücksempfinden in
insgesamt 165 Ländern die über die Jahre 2005 bis 2018 erhoben wurden.
Es gibt 24 Variablen die für ein Land in einem Jahr jeweils erhoben
wurden. In dem vorliegenden Datensatz liegen insgesamt 1784
Beobachtungen vor.

Um zu ermitteln welche Länder den höchsten Zufriedenheitswert haben wird
der jeweilige Durchschnitt über die in den Daten vorhandenen Jahre
beobachtete Zufriedenheit berechnet. Anschließend werden die Länder nach
diesem Durchschnittswert sortiert. Woraus sich eine Liste ergibt welche
Länder sich gut und nicht gut zum auswandern eignen.

\begin{Shaded}
\begin{Highlighting}[]
\CommentTok{\# Grafiken im Mittel}
\FunctionTok{ggplot}\NormalTok{(best\_worst, }\FunctionTok{aes}\NormalTok{(}\AttributeTok{x =} \FunctionTok{reorder}\NormalTok{(Country\_name, Life\_Ladder),}
                       \AttributeTok{y =}\NormalTok{ Life\_Ladder,}
                       \AttributeTok{fill =}\NormalTok{ Life\_Ladder)) }\SpecialCharTok{+}
  \FunctionTok{geom\_bar}\NormalTok{(}\AttributeTok{stat =} \StringTok{"identity"}\NormalTok{) }\SpecialCharTok{+}
  \FunctionTok{scale\_fill\_gradient}\NormalTok{(}\AttributeTok{low =} \StringTok{"tomato"}\NormalTok{, }\AttributeTok{high =} \StringTok{"darkgreen"}\NormalTok{, }\AttributeTok{guide =} \StringTok{"none"}\NormalTok{) }\SpecialCharTok{+}  
  \FunctionTok{geom\_text}\NormalTok{(}\FunctionTok{aes}\NormalTok{(}\AttributeTok{label =} \FunctionTok{round}\NormalTok{(Life\_Ladder, }\DecValTok{2}\NormalTok{)),                            }
            \AttributeTok{vjust =} \SpecialCharTok{{-}}\FloatTok{0.3}\NormalTok{, }\AttributeTok{size =} \FloatTok{3.5}\NormalTok{) }\SpecialCharTok{+}
  \FunctionTok{labs}\NormalTok{(}\AttributeTok{x =} \StringTok{"Countries"}\NormalTok{,}
       \AttributeTok{y =} \StringTok{"Life Ladder Score [mean]"}\NormalTok{,}
       \AttributeTok{title =} \StringTok{"Worst/Best five Countries in Happiness Score"}\NormalTok{) }\SpecialCharTok{+}
  \FunctionTok{theme\_minimal}\NormalTok{() }\SpecialCharTok{+}
  \FunctionTok{theme}\NormalTok{(}\AttributeTok{axis.text.x =} \FunctionTok{element\_text}\NormalTok{(}\AttributeTok{angle =} \DecValTok{45}\NormalTok{, }\AttributeTok{hjust =} \DecValTok{1}\NormalTok{)) }\SpecialCharTok{+}
  \FunctionTok{coord\_cartesian}\NormalTok{(}\AttributeTok{ylim =} \FunctionTok{c}\NormalTok{(}\DecValTok{0}\NormalTok{,}\FunctionTok{max}\NormalTok{(best\_worst}\SpecialCharTok{$}\NormalTok{Life\_Ladder)}\SpecialCharTok{*}\FloatTok{1.1}\NormalTok{))}
\end{Highlighting}
\end{Shaded}

\includegraphics{Inspection_files/figure-latex/unnamed-chunk-3-1.pdf}

\begin{Shaded}
\begin{Highlighting}[]
\CommentTok{\# Aufgabe 1)}
\CommentTok{\# Gesamtscore:}
\CommentTok{\# Welche Länder haben den h¨ochsten/niedrigsten Gesamtscore und eignen sich deswegen besonders}
\CommentTok{\# als neue Heimat bzw. sind eher auszuschließen? Welche Faktoren haben die st¨arkste Verbindung}
\CommentTok{\# mit dem Gesamtscore?}
\end{Highlighting}
\end{Shaded}

In der Abbildung sind die fünf Länder mit dem durchschnittlich besten
und schlechtesten Wert abgebildet. Die Länder welche den höchsten
durchschnittlichen Wert haben und sich damit am besten zum auswandern
eignen sind die Niederlande mit einem durchschnittlichen Wert von
ungefähr 7.47, die Schweiz mit einem Durchschnitt von 7.53, Norwegen und
Finnland mit einem Wert von je 7.55 und Dänemark mit dem höchsten
durchschnittlichen Wert von 7.69. Zu den Ländern mit dem schlechtesten
durchschnittlichen Wert zählen Süd Sudan, Togo, die Zentralafrikanische
Republik, Burundi sowie Ruanda. Diese Reihenfolge basiert auf dem
arithmetischen Mittel, dieses hat einen Bruchpunkt von 1 und reagiert
damit im allgemeinen sensibel auf Ausreißer. Deswegen wird zum einen der
Median bestimmt um die Robustheit dieser Ergebnisse zu überprüfen.

\begin{Shaded}
\begin{Highlighting}[]
\FunctionTok{ggplot}\NormalTok{(best\_worst\_med, }\FunctionTok{aes}\NormalTok{(}\AttributeTok{x =} \FunctionTok{reorder}\NormalTok{(Country\_name, Life\_Ladder),}
                       \AttributeTok{y =}\NormalTok{ Life\_Ladder,}
                       \AttributeTok{fill =}\NormalTok{ Life\_Ladder)) }\SpecialCharTok{+}
  \FunctionTok{geom\_bar}\NormalTok{(}\AttributeTok{stat =} \StringTok{"identity"}\NormalTok{) }\SpecialCharTok{+}
  \FunctionTok{scale\_fill\_gradient}\NormalTok{(}\AttributeTok{low =} \StringTok{"tomato"}\NormalTok{, }\AttributeTok{high =} \StringTok{"darkgreen"}\NormalTok{, }\AttributeTok{guide =} \StringTok{"none"}\NormalTok{) }\SpecialCharTok{+}  
  \FunctionTok{geom\_text}\NormalTok{(}\FunctionTok{aes}\NormalTok{(}\AttributeTok{label =} \FunctionTok{round}\NormalTok{(Life\_Ladder, }\DecValTok{2}\NormalTok{)),                            }
            \AttributeTok{vjust =} \SpecialCharTok{{-}}\FloatTok{0.3}\NormalTok{, }\AttributeTok{size =} \FloatTok{3.5}\NormalTok{) }\SpecialCharTok{+}
  \FunctionTok{labs}\NormalTok{(}\AttributeTok{x =} \StringTok{"Countries"}\NormalTok{,}
       \AttributeTok{y =} \StringTok{"Life Ladder Score [median]"}\NormalTok{,}
       \AttributeTok{title =} \StringTok{"Worst/Best five Countries in Happiness Score"}\NormalTok{) }\SpecialCharTok{+}
  \FunctionTok{theme\_minimal}\NormalTok{() }\SpecialCharTok{+}
  \FunctionTok{theme}\NormalTok{(}\AttributeTok{axis.text.x =} \FunctionTok{element\_text}\NormalTok{(}\AttributeTok{angle =} \DecValTok{45}\NormalTok{, }\AttributeTok{hjust =} \DecValTok{1}\NormalTok{)) }\SpecialCharTok{+}
  \FunctionTok{coord\_cartesian}\NormalTok{(}\AttributeTok{ylim =} \FunctionTok{c}\NormalTok{(}\DecValTok{0}\NormalTok{,}\FunctionTok{max}\NormalTok{(best\_worst}\SpecialCharTok{$}\NormalTok{Life\_Ladder)}\SpecialCharTok{*}\FloatTok{1.1}\NormalTok{))}
\end{Highlighting}
\end{Shaded}

\includegraphics{Inspection_files/figure-latex/unnamed-chunk-6-1.pdf}

Diese Abbildung bestätigt die Ergebnisse aus der Abbildung die auf den
Werten basierend auf dem arithmetischen Mittel basierten. Die fünf
Länder mit dem schlechtesten Wert (im Median) sind bis auf Burundi,
stattdessen gehört Benin dazu identisch. Bei den fünf Ländern mit dem
besten Wert (im Median) gehört nun Island statt Finnland zu den fünf
Ländern mit dem höchsten Wert. Dänemark ist auch wenn man den Median als
Metrik verwendet den höchsten Wert und auch die Reihenfolge der übrigen
Länder weist eine Ähnlichkeit zu der vorherigen Abbildung auf.

Außerdem werden die einzelnen Jahre für die Daten existieren in einem
zeitlichen Verlauf dargestellt um die Streuung der Daten zu beurteilen,
dafür werden die 10 Länder aus der Grafik mit dem arithmetischen Mittel
verwendet.

\begin{Shaded}
\begin{Highlighting}[]
\FunctionTok{ggplot}\NormalTok{(}\AttributeTok{data =}\NormalTok{ time, }\FunctionTok{aes}\NormalTok{(}\AttributeTok{x =}\NormalTok{ Year, }\AttributeTok{y =}\NormalTok{ Life\_Ladder, }\AttributeTok{group =}\NormalTok{ Country\_name, }\AttributeTok{color =}\NormalTok{ col))}\SpecialCharTok{+}\FunctionTok{geom\_line}\NormalTok{()}\SpecialCharTok{+} \FunctionTok{geom\_point}\NormalTok{()}\SpecialCharTok{+}\FunctionTok{theme\_minimal}\NormalTok{()}\SpecialCharTok{+}
  \FunctionTok{labs}\NormalTok{(}\AttributeTok{x =} \StringTok{"Time Span"}\NormalTok{, }\AttributeTok{y =} \StringTok{"Life Ladder Score"}\NormalTok{, }\AttributeTok{title =} \StringTok{"Chronological Sequence of Ladder Score"}\NormalTok{)}\SpecialCharTok{+}
  \CommentTok{\# scale\_x\_continuous(breaks = sort(unique(time$Year)))+}
  \FunctionTok{scale\_color\_manual}\NormalTok{(}\AttributeTok{values =} \FunctionTok{c}\NormalTok{(}\StringTok{"Good\_Score"} \OtherTok{=} \StringTok{"darkgreen"}\NormalTok{, }\StringTok{"Bad\_Score"} \OtherTok{=} \StringTok{"tomato"}\NormalTok{))}\SpecialCharTok{+}
  \FunctionTok{facet\_wrap}\NormalTok{(}\SpecialCharTok{\textasciitilde{}}\NormalTok{Country\_name, }\AttributeTok{scales =} \StringTok{"fixed"}\NormalTok{, }\AttributeTok{ncol =} \DecValTok{5}\NormalTok{)}
\end{Highlighting}
\end{Shaded}

\includegraphics{Inspection_files/figure-latex/unnamed-chunk-8-1.pdf}

In dieser Abbildung fällt auf das die fünf Länder die im Durchschnitt
den schlechtesten Wert für Zufriedenheit haben eine größere Streuung
besitzen als die Länder die einen guten Wert für Zufriedenheit haben.
Dies gilt zum einen über die jeweiligen Zeitspannen für die Daten
existieren, als auch innerhalb der Gruppen der Länder mit dem
höchsten/niedrigsten durchschnittlichen Wert.

\begin{Shaded}
\begin{Highlighting}[]
\CommentTok{\# Abweichungen innerhalb der guten Länder}
\NormalTok{good\_sd }\OtherTok{\textless{}{-}} \FunctionTok{sd}\NormalTok{(df[df}\SpecialCharTok{$}\NormalTok{Country\_name }\SpecialCharTok{\%in\%} \FunctionTok{c}\NormalTok{(}\StringTok{"Netherlands"}\NormalTok{, }\StringTok{"Switzerland"}\NormalTok{, }\StringTok{"Norway"}\NormalTok{, }\StringTok{"Finland"}\NormalTok{, }\StringTok{"Denmark"}\NormalTok{),]}\SpecialCharTok{$}\NormalTok{Life\_Ladder)}

\CommentTok{\# Abweichung innerhalb der schlechten Länder}
\NormalTok{bad\_sd }\OtherTok{\textless{}{-}} \FunctionTok{sd}\NormalTok{(df[df}\SpecialCharTok{$}\NormalTok{Country\_name }\SpecialCharTok{\%in\%} \FunctionTok{c}\NormalTok{(}\StringTok{"South Sudan"}\NormalTok{, }\StringTok{"Togo"}\NormalTok{, }\StringTok{"Central African Republic"}\NormalTok{, }\StringTok{"Burundi"}\NormalTok{, }\StringTok{"Rwanda"}\NormalTok{),]}\SpecialCharTok{$}\NormalTok{Life\_Ladder)}

\FunctionTok{paste}\NormalTok{(}\StringTok{"Die Standardabweichung innerhalb der Länder mit dem schlechtesten Score beträgt:"}\NormalTok{,}\FunctionTok{round}\NormalTok{(bad\_sd,}\DecValTok{3}\NormalTok{), }\StringTok{"Die Standardabweichung innerhalb der Länder mit dem besten Score beträgt:"}\NormalTok{, }\FunctionTok{round}\NormalTok{(good\_sd,}\DecValTok{3}\NormalTok{))}
\end{Highlighting}
\end{Shaded}

\begin{verbatim}
## [1] "Die Standardabweichung innerhalb der Länder mit dem schlechtesten Score beträgt: 0.485 Die Standardabweichung innerhalb der Länder mit dem besten Score beträgt: 0.156"
\end{verbatim}

Es soll der Zusammenhang zwischen der ``Hauptvariable'' Life Ladder und
den restlichen erhobenen Variablen ermittelt werden. Um den Zusammenhang
zwischen Variablen zu messen benötigt man ein geeignetes
Zusammenhangsmaß.Ein solches Maß ist beispielsweise der Bravais Pearson
Korrelationskoeffizient, welcher Werte innerhalb von {[}-1,1{]} annehmen
kann, wobei ein Wert von -1 eine starke negative Korrelation ausdrückt,
ein Wert von 1 eine starke positive Korrelation beschreibt und ein Wert
von 0 unkorreliertheit zwischen zwei Variablen bedeutet.

\begin{Shaded}
\begin{Highlighting}[]
\CommentTok{\# Barplot}
\FunctionTok{ggplot}\NormalTok{(cor\_df, }\FunctionTok{aes}\NormalTok{(}\AttributeTok{x =} \FunctionTok{reorder}\NormalTok{(Variable, Correlation), }\AttributeTok{y =}\NormalTok{ Correlation, }\AttributeTok{fill =}\NormalTok{ Correlation)) }\SpecialCharTok{+}
  \FunctionTok{geom\_bar}\NormalTok{(}\AttributeTok{stat =} \StringTok{"identity"}\NormalTok{, }\AttributeTok{na.rm =} \ConstantTok{TRUE}\NormalTok{) }\SpecialCharTok{+}
  \FunctionTok{coord\_flip}\NormalTok{() }\SpecialCharTok{+}
  \FunctionTok{scale\_fill\_gradient2}\NormalTok{(}\AttributeTok{low =} \StringTok{"blue"}\NormalTok{, }\AttributeTok{mid =} \StringTok{"white"}\NormalTok{, }\AttributeTok{high =} \StringTok{"red"}\NormalTok{, }\AttributeTok{midpoint =} \DecValTok{0}\NormalTok{) }\SpecialCharTok{+}
  \FunctionTok{labs}\NormalTok{(}\AttributeTok{title =} \StringTok{"Variable Correlation with Life Ladder Score"}\NormalTok{,}
       \AttributeTok{x =} \StringTok{"Variable"}\NormalTok{, }\AttributeTok{y =} \StringTok{"Correlation"}\NormalTok{) }\SpecialCharTok{+}
  \FunctionTok{theme\_minimal}\NormalTok{() }\SpecialCharTok{+}
  \FunctionTok{theme}\NormalTok{(}\AttributeTok{axis.text =} \FunctionTok{element\_text}\NormalTok{(}\AttributeTok{size =} \DecValTok{10}\NormalTok{),}
        \AttributeTok{plot.title =} \FunctionTok{element\_text}\NormalTok{(}\AttributeTok{hjust =} \FloatTok{0.5}\NormalTok{))}
\end{Highlighting}
\end{Shaded}

\includegraphics{Inspection_files/figure-latex/unnamed-chunk-11-1.pdf}

In dieser Abbildung sind die Korrelationen jeder Variable mit Life
Ladder dargestellt. Die Korrelation von Life Ladder zu sich selber
beträgt 1 und bezeichnet eine ``perfekte'' Korrelation und ist zu
Anschauungszwecken ebenfalls im Plot enthalten. In der Statistik gibt es
den Begriff der Multikollinearität. Wenn zwei Variablen durch einen
gemeinsamen Effekt beeinflusst werden, dann werden sie obwohl sie eine
unterschiedliche Bedeutung haben eine hohe Korrelation aufweisen. Da man
dann unter Umständen diesen Effekt in mehreren Variablen merkt, ist es
sinnvoll die Kovarianzstruktur genauer zu betrachten, um nur die
Variablen zu beachten die einen hohen Beitrag zur Gesamtvarianz haben.

Um den Einfluss von Spalten mit einem hohen Anteil fehlender Werte zu
vermeiden, wurden nur jene Variablen beibehalten, für die mindestens 80
\% der Werte vorhanden waren. Die Variablen wurden dabei standardisiert
(z-Transformation), um unterschiedliche Skalierungen zu neutralisieren.
Aus der Rotationsmatrix der PCA (\$rotation) wurde der Gesamtbeitrag
jeder Variablen zu den ersten drei Hauptkomponenten berechnet. Dazu
wurde für jede Variable die Summe der quadrierten Ladungen (loadings)
über die ersten drei Komponenten gebildet. Dieses Maß beschreibt, in
welchem Ausmaß eine Variable zur erklärten Varianz der wichtigsten
Komponenten beiträgt. Auf dieser Grundlage wurden die zehn Variablen mit
dem höchsten Gesamtbeitrag ausgewählt, da sie die größte strukturelle
Bedeutung innerhalb der Daten aufweisen. Im Anschluss wurde für diese
Variablen die Pearson-Korrelation mit der Zielvariablen Life Ladder
berechnet, um zu untersuchen, welche dieser zentralen Variablen am
stärksten mit dem subjektiven Wohlbefinden (Life Ladder) zusammenhängen.
Die Hauptkomponentenanalyse wurde unabhängig von der Zielvariable Life
Ladder durchgeführt und dient nicht zur Analyse der Korrelation, sondern
als Maßnahme zur vorverarbeitung der Daten die die Dimension der Daten
verringern soll.

\begin{Shaded}
\begin{Highlighting}[]
\FunctionTok{ggplot}\NormalTok{(cor\_df[cor\_df}\SpecialCharTok{$}\NormalTok{Variable }\SpecialCharTok{\%in\%}\NormalTok{ top15,], }\FunctionTok{aes}\NormalTok{(}\AttributeTok{x =} \FunctionTok{reorder}\NormalTok{(Variable, Correlation), }\AttributeTok{y =}\NormalTok{ Correlation, }\AttributeTok{fill =}\NormalTok{ Correlation)) }\SpecialCharTok{+}
  \FunctionTok{geom\_bar}\NormalTok{(}\AttributeTok{stat =} \StringTok{"identity"}\NormalTok{, }\AttributeTok{na.rm =} \ConstantTok{TRUE}\NormalTok{) }\SpecialCharTok{+}
  \FunctionTok{coord\_flip}\NormalTok{() }\SpecialCharTok{+}
  \FunctionTok{scale\_fill\_gradient2}\NormalTok{(}\AttributeTok{low =} \StringTok{"blue"}\NormalTok{, }\AttributeTok{mid =} \StringTok{"white"}\NormalTok{, }\AttributeTok{high =} \StringTok{"red"}\NormalTok{, }\AttributeTok{midpoint =} \DecValTok{0}\NormalTok{) }\SpecialCharTok{+}
  \FunctionTok{labs}\NormalTok{(}\AttributeTok{title =} \StringTok{"Correlation mit Life\_Ladder [only 15 most important]"}\NormalTok{,}
       \AttributeTok{x =} \StringTok{"Variable"}\NormalTok{, }\AttributeTok{y =} \StringTok{"Correlation"}\NormalTok{) }\SpecialCharTok{+}
  \FunctionTok{theme\_minimal}\NormalTok{() }\SpecialCharTok{+}
  \FunctionTok{theme}\NormalTok{(}\AttributeTok{axis.text =} \FunctionTok{element\_text}\NormalTok{(}\AttributeTok{size =} \DecValTok{10}\NormalTok{),}
        \AttributeTok{plot.title =} \FunctionTok{element\_text}\NormalTok{(}\AttributeTok{hjust =} \FloatTok{0.5}\NormalTok{))}
\end{Highlighting}
\end{Shaded}

\includegraphics{Inspection_files/figure-latex/unnamed-chunk-13-1.pdf}

\begin{Shaded}
\begin{Highlighting}[]
\FunctionTok{sum}\NormalTok{(var\_df[var\_df}\SpecialCharTok{$}\NormalTok{Spalte }\SpecialCharTok{\%in\%}\NormalTok{ top15,]}\SpecialCharTok{$}\NormalTok{Varianz\_erklärt)}
\end{Highlighting}
\end{Shaded}

\begin{verbatim}
## [1] 0.9375
\end{verbatim}

Die 15 in der Grafik dargestellten Variablen weisen den höchsten Beitrag
zur Varianz auf. Insgesamt erklären machen diese 15 Spalten 93.75\% der
gesamten Varianz aus. Die stärkste positive Korrelation liegt bei dem
logarithmierten Pro-Kopf Bruttoninlandsprodukt und Life Ladder vor.

Der Gini-Koeffizient ist ein Maß zur Beschreibung der Ungleichheit einer
Verteilung. Er wird häufig verwendet, um die Verteilung von Einkommen
oder Vermögen in einer Bevölkerung zu quantifizieren. Der Wert liegt
zwischen 0 und 1: Ein Gini-Koeffizient von 0 bedeutet vollständige
Gleichheit (alle besitzen gleich viel), während ein Wert von 1 maximale
Ungleichheit beschreibt (eine Person besitzt alles, alle anderen
nichts). In der folgenden Analyse wird der Gini Koeffizient für den Life
Ladder Score, sowie für das Pro-Kopf Bruttoinlandsprodukt berechnet.

\begin{Shaded}
\begin{Highlighting}[]
\CommentTok{\# Aufgabe 2)}
\CommentTok{\# Gini Index:}
\CommentTok{\# Bestimmen Sie den Gini{-}Index aller Länder für das letzte Jahr.}
\end{Highlighting}
\end{Shaded}

\begin{Shaded}
\begin{Highlighting}[]
\CommentTok{\#Gini Koeffizient }

\CommentTok{\# last\_year \textless{}{-} df \%\textgreater{}\% group\_by("Country\_Name") \%\textgreater{}\% }


\NormalTok{means}
\end{Highlighting}
\end{Shaded}

\begin{verbatim}
## # A tibble: 165 x 25
##    Country_name Life_Ladder Log_GDP_per_capita Social_support
##    <chr>              <dbl>              <dbl>          <dbl>
##  1 Afghanistan         3.71               7.44          0.516
##  2 Albania             4.99               9.26          0.720
##  3 Algeria             5.48               9.51          0.804
##  4 Angola              4.42               8.71          0.738
##  5 Argentina           6.36               9.83          0.906
##  6 Armenia             4.44               8.93          0.714
##  7 Australia           7.30              10.7           0.949
##  8 Austria             7.25              10.7           0.927
##  9 Azerbaijan          4.92               9.63          0.762
## 10 Bahrain             5.86              10.6           0.884
## # i 155 more rows
## # i 21 more variables: Healthy_life_expectancy_at_birth <dbl>,
## #   Freedom_to_make_life_choices <dbl>, Generosity <dbl>,
## #   Perceptions_of_corruption <dbl>, Positive_affect <dbl>,
## #   Negative_affect <dbl>, Confidence_in_national_government <dbl>,
## #   Democratic_Quality <dbl>, Delivery_Quality <dbl>,
## #   `Standard_deviation_of_ladder_by_country-year` <dbl>, ...
\end{verbatim}

\begin{Shaded}
\begin{Highlighting}[]
\NormalTok{means }\SpecialCharTok{\%\textgreater{}\%} \FunctionTok{summarise}\NormalTok{(}\AttributeTok{Gini\_Score\_LGDP =} \FunctionTok{Gini}\NormalTok{(}\FunctionTok{exp}\NormalTok{(Log\_GDP\_per\_capita), }\AttributeTok{na.rm =} \ConstantTok{TRUE}\NormalTok{))}
\end{Highlighting}
\end{Shaded}

\begin{verbatim}
## # A tibble: 1 x 1
##   Gini_Score_LGDP
##             <dbl>
## 1           0.535
\end{verbatim}

\begin{Shaded}
\begin{Highlighting}[]
\NormalTok{means }\SpecialCharTok{\%\textgreater{}\%} \FunctionTok{summarise}\NormalTok{(}\AttributeTok{Gini\_Score\_Ladder =} \FunctionTok{Gini}\NormalTok{(Life\_Ladder, }\AttributeTok{na.rm =} \ConstantTok{TRUE}\NormalTok{))}
\end{Highlighting}
\end{Shaded}

\begin{verbatim}
## # A tibble: 1 x 1
##   Gini_Score_Ladder
##               <dbl>
## 1             0.113
\end{verbatim}

\begin{Shaded}
\begin{Highlighting}[]
\NormalTok{df }\SpecialCharTok{\%\textgreater{}\%} \FunctionTok{summarise}\NormalTok{(}\AttributeTok{Gini\_Score\_LGDP =} \FunctionTok{Gini}\NormalTok{(}\FunctionTok{exp}\NormalTok{(Log\_GDP\_per\_capita), }\AttributeTok{na.rm =} \ConstantTok{TRUE}\NormalTok{))}
\end{Highlighting}
\end{Shaded}

\begin{verbatim}
##   Gini_Score_LGDP
## 1       0.5173427
\end{verbatim}

\begin{Shaded}
\begin{Highlighting}[]
\NormalTok{df }\SpecialCharTok{\%\textgreater{}\%} \FunctionTok{filter}\NormalTok{(Year }\SpecialCharTok{==} \DecValTok{2018}\NormalTok{) }\SpecialCharTok{\%\textgreater{}\%} \FunctionTok{summarise}\NormalTok{(}\AttributeTok{Gini\_Score\_Ladder =} \FunctionTok{Gini}\NormalTok{(Life\_Ladder, }\AttributeTok{na.rm =} \ConstantTok{TRUE}\NormalTok{))}
\end{Highlighting}
\end{Shaded}

\begin{verbatim}
##   Gini_Score_Ladder
## 1         0.1139761
\end{verbatim}

\begin{Shaded}
\begin{Highlighting}[]
\CommentTok{\# Unterschied zwischen der Verteilung von Life Ladder und von Log GDP per Capita}


\FunctionTok{ggplot}\NormalTok{(means, }\FunctionTok{aes}\NormalTok{(}\AttributeTok{x =}\NormalTok{ (Log\_GDP\_per\_capita), }\AttributeTok{y =}\NormalTok{ Life\_Ladder)) }\SpecialCharTok{+}
  \FunctionTok{geom\_point}\NormalTok{(}\FunctionTok{aes}\NormalTok{(}\AttributeTok{color =}\NormalTok{ Country\_name), }\AttributeTok{alpha =} \FloatTok{0.7}\NormalTok{) }\SpecialCharTok{+}
  \FunctionTok{geom\_smooth}\NormalTok{(}\AttributeTok{method =} \StringTok{"lm"}\NormalTok{, }\AttributeTok{se =} \ConstantTok{FALSE}\NormalTok{, }\AttributeTok{color =} \StringTok{"black"}\NormalTok{) }\SpecialCharTok{+}
  \FunctionTok{theme\_minimal}\NormalTok{() }\SpecialCharTok{+}
  \FunctionTok{labs}\NormalTok{(}
    \AttributeTok{title =} \StringTok{"Zusammenhang zwischen Log GDP per Capita und Life Ladder"}\NormalTok{,}
    \AttributeTok{x =} \StringTok{"Durchschnittliches Log GDP per Capita"}\NormalTok{,}
    \AttributeTok{y =} \StringTok{"Durchschnittlicher Life Ladder Score"}\NormalTok{,}
    \AttributeTok{color =} \StringTok{"Land"}
\NormalTok{  ) }\SpecialCharTok{+} \FunctionTok{theme}\NormalTok{(}\AttributeTok{legend.position =} \StringTok{"none"}\NormalTok{)}
\end{Highlighting}
\end{Shaded}

\begin{verbatim}
## `geom_smooth()` using formula = 'y ~ x'
\end{verbatim}

\begin{verbatim}
## Warning: Removed 3 rows containing non-finite outside the scale range
## (`stat_smooth()`).
\end{verbatim}

\begin{verbatim}
## Warning: Removed 3 rows containing missing values or values outside the scale range
## (`geom_point()`).
\end{verbatim}

\includegraphics{Inspection_files/figure-latex/unnamed-chunk-16-1.pdf}

\begin{Shaded}
\begin{Highlighting}[]
\NormalTok{df\_sorted\_GDP }\OtherTok{\textless{}{-}}\NormalTok{ means[}\FunctionTok{order}\NormalTok{(means}\SpecialCharTok{$}\NormalTok{Log\_GDP\_per\_capita), ]}

\NormalTok{df\_sorted\_GDP}\SpecialCharTok{$}\NormalTok{cum\_share }\OtherTok{\textless{}{-}} \FunctionTok{cumsum}\NormalTok{(}\FunctionTok{exp}\NormalTok{(df\_sorted\_GDP}\SpecialCharTok{$}\NormalTok{Log\_GDP\_per\_capita)) }\SpecialCharTok{/} \FunctionTok{sum}\NormalTok{(}\FunctionTok{exp}\NormalTok{(df\_sorted\_GDP}\SpecialCharTok{$}\NormalTok{Log\_GDP\_per\_capita), }\AttributeTok{na.rm =} \ConstantTok{TRUE}\NormalTok{)}


\FunctionTok{length}\NormalTok{(}\FunctionTok{na.omit}\NormalTok{(df\_sorted\_GDP[df\_sorted\_GDP}\SpecialCharTok{$}\NormalTok{cum\_share }\SpecialCharTok{\textgreater{}} \FloatTok{0.5}\NormalTok{,]}\SpecialCharTok{$}\NormalTok{Country\_name))}
\end{Highlighting}
\end{Shaded}

\begin{verbatim}
## [1] 28
\end{verbatim}

\begin{Shaded}
\begin{Highlighting}[]
\NormalTok{df\_sorted\_LL }\OtherTok{\textless{}{-}}\NormalTok{ means[}\FunctionTok{order}\NormalTok{(means}\SpecialCharTok{$}\NormalTok{Life\_Ladder), ]}

\NormalTok{df\_sorted\_LL}\SpecialCharTok{$}\NormalTok{cum\_share }\OtherTok{\textless{}{-}} \FunctionTok{cumsum}\NormalTok{(df\_sorted\_LL}\SpecialCharTok{$}\NormalTok{Life\_Ladder) }\SpecialCharTok{/} \FunctionTok{sum}\NormalTok{(df\_sorted\_GDP}\SpecialCharTok{$}\NormalTok{Life\_Ladder, }\AttributeTok{na.rm =} \ConstantTok{TRUE}\NormalTok{)}

\FunctionTok{length}\NormalTok{(}\FunctionTok{na.omit}\NormalTok{(df\_sorted\_LL[df\_sorted\_LL}\SpecialCharTok{$}\NormalTok{cum\_share }\SpecialCharTok{\textgreater{}} \FloatTok{0.5}\NormalTok{,]}\SpecialCharTok{$}\NormalTok{Country\_name))}
\end{Highlighting}
\end{Shaded}

\begin{verbatim}
## [1] 70
\end{verbatim}

\begin{Shaded}
\begin{Highlighting}[]
\CommentTok{\# Lebenszufriedenheit ist nicht so ungleich verteilt wie das Pro Kopf Einkommen}
\end{Highlighting}
\end{Shaded}

Bei der Betrachtung der beiden Gini Koeffizienten fällt auf das Pro-Kopf
Bruttoinlandsprodukt stärker auf wenige Länder konzentriert ist als der
Life Ladder Score.

\begin{Shaded}
\begin{Highlighting}[]
\FunctionTok{summary}\NormalTok{(}\FunctionTok{lm}\NormalTok{(}\AttributeTok{data =}\NormalTok{ means,}\AttributeTok{formula =}\NormalTok{ Life\_Ladder }\SpecialCharTok{\textasciitilde{}}\NormalTok{ Log\_GDP\_per\_capita }\SpecialCharTok{+}\NormalTok{ Social\_support }\SpecialCharTok{+}\NormalTok{ Democratic\_Quality))}
\end{Highlighting}
\end{Shaded}

\begin{verbatim}
## 
## Call:
## lm(formula = Life_Ladder ~ Log_GDP_per_capita + Social_support + 
##     Democratic_Quality, data = means)
## 
## Residuals:
##      Min       1Q   Median       3Q      Max 
## -1.79632 -0.32951 -0.00002  0.36076  1.29163 
## 
## Coefficients:
##                    Estimate Std. Error t value Pr(>|t|)    
## (Intercept)        -0.79003    0.44021  -1.795  0.07463 .  
## Log_GDP_per_capita  0.41814    0.05478   7.633 2.08e-12 ***
## Social_support      2.94859    0.54230   5.437 2.03e-07 ***
## Democratic_Quality  0.22597    0.06325   3.572  0.00047 ***
## ---
## Signif. codes:  0 '***' 0.001 '**' 0.01 '*' 0.05 '.' 0.1 ' ' 1
## 
## Residual standard error: 0.5297 on 157 degrees of freedom
##   (4 Beobachtungen als fehlend gelöscht)
## Multiple R-squared:  0.7612, Adjusted R-squared:  0.7566 
## F-statistic: 166.8 on 3 and 157 DF,  p-value: < 2.2e-16
\end{verbatim}

\begin{Shaded}
\begin{Highlighting}[]
\FunctionTok{summary}\NormalTok{(}\FunctionTok{lm}\NormalTok{(}\AttributeTok{data =}\NormalTok{ means, }\AttributeTok{formula =}\NormalTok{ Life\_Ladder }\SpecialCharTok{\textasciitilde{}}\NormalTok{ Healthy\_life\_expectancy\_at\_birth))}
\end{Highlighting}
\end{Shaded}

\begin{verbatim}
## 
## Call:
## lm(formula = Life_Ladder ~ Healthy_life_expectancy_at_birth, 
##     data = means)
## 
## Residuals:
##      Min       1Q   Median       3Q      Max 
## -1.50561 -0.49057 -0.00576  0.52262  1.36856 
## 
## Coefficients:
##                                   Estimate Std. Error t value Pr(>|t|)    
## (Intercept)                      -1.586377   0.427370  -3.712 0.000284 ***
## Healthy_life_expectancy_at_birth  0.111652   0.006791  16.441  < 2e-16 ***
## ---
## Signif. codes:  0 '***' 0.001 '**' 0.01 '*' 0.05 '.' 0.1 ' ' 1
## 
## Residual standard error: 0.6573 on 160 degrees of freedom
##   (3 Beobachtungen als fehlend gelöscht)
## Multiple R-squared:  0.6282, Adjusted R-squared:  0.6258 
## F-statistic: 270.3 on 1 and 160 DF,  p-value: < 2.2e-16
\end{verbatim}

\begin{Shaded}
\begin{Highlighting}[]
\CommentTok{\# Fehlende Werte:}
\CommentTok{\# In einigen Bereichen ist der Datensatz nicht komplett. Wie kann das zustande kommen und wie}
\CommentTok{\# könnte man bei der Deskription darauf eingehen?}
\end{Highlighting}
\end{Shaded}

\begin{Shaded}
\begin{Highlighting}[]
\CommentTok{\# Da Paul gerne asiatisch isst, interessiert ihn besonders der Vergleich von Hong Kong, China und}
\CommentTok{\# Deutschland. Dieser sollte f¨ur ihn entsprechend optisch aufbereitet werden.}
\end{Highlighting}
\end{Shaded}

\begin{Shaded}
\begin{Highlighting}[]
\CommentTok{\# Untersuchen Sie zus¨atzlich 3 weitere Aspekte Ihrer Wahl.}
\end{Highlighting}
\end{Shaded}


\end{document}
